\chapter{C++ for Latency II: Compile Time}\label{ch:ll:cpp-for-latency-ii-compile-time}

Four stages of a hot path, each a class with a virtual function called through a pointer, cost this book's laptop about
seven nanoseconds per event. The same four stages composed as template parameters cost about one: the compiler, seeing
every call, inlined all four into thirteen machine instructions with no call left in them. The work was identical; what
changed is how much the compiler knew when it compiled. This chapter is about moving decisions from run time to compile
time in C++: dispatch by templates instead of virtual functions, values computed by the compiler, \omterm{def:ll:cpp-for-latency-ii-compile-time:inline}{inlining} and what
prevents it, and the layout of rare code away from the common path.

\section{Virtual dispatch and what it costs}

\begin{definition}[Dynamic dispatch, devirtualisation]\label{def:ll:cpp-for-latency-ii-compile-time:dyn}
In \emph{dynamic dispatch}\index{dynamic dispatch} the function called depends on the run-time type of an object: a C++
virtual call loads the object's table of function pointers and calls through it, an indirect call. \emph{Devirtualisation}\index{devirtualisation}
is the compiler's replacement of such a call by a direct one, when it can prove the object's type (a \texttt{final}
class, a local object, whole-program analysis).
\end{definition}

An indirect call costs little when the processor predicts its target: the branch predictor keeps a history of targets,
and a hot path that always sees the same message type always calls the same function. It costs a misprediction when the
target changes unpredictably, and it always costs the optimisations it blocks: the compiler cannot inline a function it
cannot name, so it cannot fold the callee's work into the caller, keep values in registers across the call, or vectorise
the loop around it. The published measurement of the direct cost of virtual calls in C++ (Driesen and Hölzle, 1996) found
it modest in cycles; the indirect cost, lost optimisation, is usually larger and does not appear in a \omterm{def:ll:cpu-microarchitecture:ubench}{microbenchmark} of the
call alone.

\section{Static polymorphism: templates, CRTP and concepts}

\begin{definition}[Static polymorphism, curiously recurring template pattern]\label{def:ll:cpp-for-latency-ii-compile-time:static}
\emph{Static polymorphism}\index{static polymorphism} selects the function to call at compile time, by templates or
overloading, so that every call is direct. In the \emph{curiously recurring template pattern}\index{curiously recurring
template pattern} (CRTP) a class template takes its own derived class as a parameter, \texttt{struct Book : Base<Book>},
and the base calls the derived class's functions through a \texttt{static\_cast} of \texttt{this}.
\end{definition}

A trading hot path rarely needs run-time polymorphism: the set of message types is fixed by the protocol, and the strategy
is chosen when the process starts. Both can be template parameters. C++20 concepts state the requirement a parameter must
meet (``has \texttt{bool on(E\&)}''), so the error for a stage that does not fit is one line instead of a page.
\Cref{fig:ll:cpp-for-latency-ii-compile-time:dispatch} compares three dispatchers of one message to one of four handlers
(add, cancel, execute, trade), each doing an addition or two: virtual handlers in a table, \texttt{std::variant} with
\texttt{std::visit}, and CRTP with a \texttt{switch}. On a stream of one message type, where every branch is predicted, the
CRTP dispatcher costs about half a nanosecond per message and the virtual one about four times as much; on a random mix of
the four types all three pay mispredictions, and the gap narrows. The last two bars are the four-stage pipeline of the
opening, static and dynamic.

\begin{omfigure}
\begin{tikzpicture}
\begin{axis}[width=12cm,height=5.6cm,ybar,bar width=10pt,ymin=0,
  symbolic x coords={virtual,variant,CRTP,pipeline static,pipeline dynamic},xtick=data,enlarge x limits=0.12,
  xticklabels={virtual,variant,CRTP,{pipeline\\ static},{pipeline\\ dynamic}},xticklabel style={align=center},
  ylabel={ns per message},tick label style={font=\small},label style={font=\small},grid=major,grid style={gray!25},
  legend columns=2,legend style={font=\footnotesize,at={(0.5,-0.33)},anchor=north,draw=none,fill=none,
  /tikz/every even column/.append style={column sep=8pt}},table/col sep=comma]
\addplot[fill=omDef!55,draw=omDef] table[x=design,y=one_type]{figdata/low-latency/07-cpp-for-latency-ii-compile-time/measured_dispatch.csv};
\addplot[fill=omThm!45,draw=omThm] table[x=design,y=mixed]{figdata/low-latency/07-cpp-for-latency-ii-compile-time/measured_dispatch.csv};
\legend{one message type,random mix of four}
\end{axis}
\end{tikzpicture}

\omcaption{Cost per message of three ways to dispatch it to its handler, on a stream of one type and on a random mix, and of
a four-stage pipeline composed statically (\texttt{firm.pipeline}) and through virtual calls (the pipeline's cost does not
depend on the message mix). Measured on a laptop (Intel Core Ultra 7 155H) under WSL2, no isolated cores, pinned to one CPU,
best of 15 runs of 65\,536 messages. Data: \texttt{bench\_dispatch.py}.}\label{fig:ll:cpp-for-latency-ii-compile-time:dispatch}
\end{omfigure}

\section{Compile-time evaluation}

\begin{definition}[Compile-time evaluation]\label{def:ll:cpp-for-latency-ii-compile-time:consteval}
\emph{Compile-time evaluation}\index{compile-time evaluation} computes a value while compiling instead of while running:
a \texttt{constexpr} function may be evaluated by the compiler when its arguments are constants, and a \texttt{consteval}
function must be, so that its result is a constant embedded in the program.
\end{definition}

Tables are the typical use: powers of ten to rescale prices between protocols with different implied decimals, the
mapping from a protocol's message type byte to a handler index, the lengths of fixed-size messages, checksum tables. Built
by the compiler, a table costs nothing at start-up, cannot be forgotten in an initialisation routine, and is checked by
\texttt{static\_assert} when the program compiles. Layout facts belong there too: the build of chapter~3 asserts at
compile time that a padded counter occupies exactly one \omterm{def:ll:memory-hierarchy-and-caches:line}{cache line}, and chapter~16 asserts each message's size against
the protocol's specification.

\section{Inlining and its limits}

\begin{definition}[Inlining]\label{def:ll:cpp-for-latency-ii-compile-time:inline}
\emph{Inlining}\index{inlining} is the compiler's replacement of a call by the body of the called function, which removes
the call's own cost and, more importantly, lets the compiler optimise the combined code: propagate constants, keep values
in registers, remove work whose result is unused.
\end{definition}

The static pipeline of \cref{lst:ll:cpp-for-latency-ii-compile-time:asm-static} is the whole effect in thirteen
instructions: the four stages' additions, the exclusive-or, the comparison and the conditional update, with the stage
objects folded into loads and stores. Its dynamic twin (\cref{lst:ll:cpp-for-latency-ii-compile-time:asm-dyn}) is a loop
over a vector of pointers with an indirect call per stage. The compiler inlines by a cost model: small functions, functions
called once, functions whose definition it can see. It will not inline across translation units without link-time
optimisation (chapter~8), through a function pointer it cannot resolve, or a function marked \texttt{noinline}; and it may
decline to inline a large function into a hot loop. \texttt{[[gnu::always\_inline]]} overrides the cost model for the few
functions that must be inlined; it is a statement about this compiler's heuristics, and should be verified in the
assembly rather than trusted.

\section{Branch layout and cold paths}

\begin{definition}[Cold path]\label{def:ll:cpp-for-latency-ii-compile-time:cold}
A \emph{cold path}\index{cold path} is code that runs rarely (error handling, logging of an anomaly, recovery). Keeping it
out of the hot function, by marking the rare branch \texttt{[[unlikely]]} and the rare function \texttt{cold} and
\texttt{noinline}, keeps the hot code small and contiguous in the instruction cache, with the common case falling through.
\end{definition}

GCC places functions marked \texttt{cold} in a separate section with the other cold functions and treats branches to them
as unlikely. The rescaling function of the tutorial checks its arguments, and the check costs one comparison and a
predicted branch; the report of a bad argument, with its formatted output, lives out of line. The discipline matters more
than any single attribute: the hot function should contain the common case and calls to everything else.

\section{Tutorial: three dispatchers and a pipeline}\label{tut:ll:cpp-for-latency-ii-compile-time:tut}

\begin{tutorial}
\textbf{Goal.} Dispatch messages three ways, compose a pipeline at compile time, and read in the assembly what the compiler
made of each. \textbf{End state:} \cref{fig:ll:cpp-for-latency-ii-compile-time:dispatch} and the two listings of assembly.
\begin{tutsteps}
\item \textbf{CRTP.} The base's \texttt{on} switches on the type and calls the derived class's handler directly.
\omcode{code/low-latency/07-cpp-for-latency-ii-compile-time/cpp/ll_dispatch.hpp}{45}{64}{Static dispatch with the curiously recurring template pattern.}\label{lst:ll:cpp-for-latency-ii-compile-time:crtp}
\item \textbf{A table the compiler builds,} and a rare path it keeps out of line.
\omcode{code/low-latency/07-cpp-for-latency-ii-compile-time/cpp/ll_dispatch.hpp}{66}{85}{A compile-time table and a cold error path.}\label{lst:ll:cpp-for-latency-ii-compile-time:consteval}
\item \textbf{Compose stages.} A fold expression over a tuple of stages, constrained by a concept; \texttt{\&\&} stops at
the first stage that rejects the event.
\omcode{code/firm/pipeline/cpp/firm_pipeline.hpp}{19}{33}{The static pipeline of firm.pipeline.}\label{lst:ll:cpp-for-latency-ii-compile-time:pipeline}
\item \textbf{Read the assembly.} \texttt{python ll\_asm.py} compiles two one-line functions with \texttt{-O2
-masm=intel} and extracts their bodies.
\omcode{code/low-latency/07-cpp-for-latency-ii-compile-time/cpp/asm/run_static.s}{1}{14}{The four-stage static pipeline after inlining (g++ 11.4, -O2).}\label{lst:ll:cpp-for-latency-ii-compile-time:asm-static}
\omcode{code/low-latency/07-cpp-for-latency-ii-compile-time/cpp/asm/run_dynamic.s}{1}{31}{The same pipeline through virtual calls: a loop with an indirect call per stage.}\label{lst:ll:cpp-for-latency-ii-compile-time:asm-dyn}
\item \textbf{Measure} with \texttt{python bench\_dispatch.py}.
\end{tutsteps}
\textbf{What to change next.} Mark the dispatch table's handlers \texttt{final} and access them through their concrete types,
and see whether the compiler devirtualises; add \texttt{[[gnu::noinline]]} to one stage of the static pipeline and read the
assembly again.
\end{tutorial}

\section{Build: the compile-time pipeline}\label{bld:ll:cpp-for-latency-ii-compile-time:pipeline}

\begin{build}
\bfield{Purpose} The composition of the hot path of Part~IV (decode, book, strategy, risk, gateway) as one statically
dispatched call chain, with a runtime-polymorphic twin for tests and for configurations chosen at start-up.
\bfield{Interface} C++20 \texttt{firm::pipeline}: concept \texttt{Stage<S, E>}; \texttt{Pipeline<E, S\dots>} with
\texttt{on(E\&)} and \texttt{stage<I>()}; \texttt{AnyStage<E>}, \texttt{DynPipeline<E>} with \texttt{add} and \texttt{on}.
Rust \texttt{firm\_pipeline}: trait \texttt{Stage<E>}, \texttt{Chain<A, B>} and the \texttt{chain!} macro,
\texttt{DynPipeline<E>}.
\bfield{Rules} Stages run in order; the first \texttt{false} stops the event; the static and dynamic pipelines give the same
results and the same order of effects; a stage that does not meet the concept does not compile.
\bfield{Acceptance tests} \texttt{code/firm/pipeline/}: order and short-circuit on a traced event, static against dynamic, in
C++20 and Rust; the chapter's test checks that the static pipeline's assembly contains no call and the dynamic one an
indirect call.
\bfield{Stretch} Stages that return a new event type (a decoded message into a book update), with the chain's types checked
by the compiler.
\end{build}

\begin{omsources}
\item K.~Driesen and U.~Hölzle, ``The direct cost of virtual function calls in C++'', OOPSLA 1996.
\item GCC manual, ``Common function attributes'' (\texttt{cold}, \texttt{noinline}, \texttt{always\_inline}).
\item cppreference, \texttt{consteval}; attributes \texttt{likely} and \texttt{unlikely}.
\item C.~Cook, ``When a microsecond is an eternity'', CppCon 2017 (templates for configuration).
\end{omsources}

\section{Exercises}

\begin{exercise}[$\star$]\label{exo:ll:cpp-for-latency-ii-compile-time:1}
From \cref{fig:ll:cpp-for-latency-ii-compile-time:dispatch}, how many times more does virtual dispatch cost than CRTP on one
message type, and on the mix?
\end{exercise}

\begin{exercise}[$\star$]\label{exo:ll:cpp-for-latency-ii-compile-time:2}
At 5 million messages a second, how much processing time a second does the dynamic pipeline cost over the static one?
\end{exercise}

\begin{exercise}[$\star$]\label{exo:ll:cpp-for-latency-ii-compile-time:3}
Write the value of \texttt{rescale(250, 2, 4)} and \texttt{rescale(2500, 4, 2)}, and say what \texttt{rescale(1, 2, 20)}
does.
\end{exercise}

\begin{exercise}[$\star\star$]\label{exo:ll:cpp-for-latency-ii-compile-time:4}
Why does the random mix slow down the CRTP dispatcher too, although it makes no indirect call?
\end{exercise}

\begin{exercise}[$\star\star$]\label{exo:ll:cpp-for-latency-ii-compile-time:5}
Count the instructions of \cref{lst:ll:cpp-for-latency-ii-compile-time:asm-static} that belong to each stage. Which stage's
object has disappeared entirely, and why?
\end{exercise}

\begin{exercise}[$\star\star$]\label{exo:ll:cpp-for-latency-ii-compile-time:6}
A strategy is chosen by a configuration file at start-up among three. How do you keep dispatch static?
\end{exercise}

\begin{exercise}[$\star\star\star$]\label{exo:ll:cpp-for-latency-ii-compile-time:7}
\emph{Coding.} Make one stage of the static pipeline \texttt{[[gnu::noinline]]}, regenerate the assembly and measure. What
changes in the listing and in the time per event?
\end{exercise}

\begin{exercise}[$\star\star\star$]\label{exo:ll:cpp-for-latency-ii-compile-time:8}
\emph{Find the flaw.} ``A virtual call costs 1.9 nanoseconds in our benchmark and our path makes five, so removing them
saves ten nanoseconds at most.''
\end{exercise}

\section{Problem: The Indirect Branch}

\begin{problem}[{Weekend problem --- a plug-in architecture on the hot path}]\label{pb:ll:cpp-for-latency-ii-compile-time:1}
A trading engine was designed with plug-ins: the decoder, book, strategy and risk check are loaded at start-up and called
through a base class. It handles four message types in a mix close to random. The team measures the alternatives of this
chapter (\cref{fig:ll:cpp-for-latency-ii-compile-time:dispatch}).

\medskip\noindent\textbf{Part I --- The measurement.}
\begin{enumerate}
\item What do the three dispatchers cost per message on one type and on the mix?
\item Which effect makes the mix expensive for all three?
\item What does the four-stage pipeline cost static and dynamic?
\item Why does the pipeline's cost not depend on the mix?
\end{enumerate}

\medskip\noindent\textbf{Part II --- The mechanisms.}
\begin{enumerate}[resume]
\item What does the dynamic pipeline do per stage that the static one does not (\cref{lst:ll:cpp-for-latency-ii-compile-time:asm-dyn})?
\item How many instructions does the static pipeline execute when the event passes every stage?
\item What optimisations became possible once the stages were inlined?
\item What would \texttt{final} on the plug-in classes change, and what not?
\end{enumerate}

\medskip\noindent\textbf{Part III --- The redesign.}
\begin{enumerate}[resume]
\item How can plug-ins chosen at start-up still be dispatched statically?
\item What is lost: what can no longer change without a rebuild?
\item How many program variants does the choice of three strategies and two risk checks produce?
\item How should the variants be tested?
\end{enumerate}

\medskip\noindent\textbf{Part IV --- The verdict.}
\begin{enumerate}[resume]
\item State the \emph{named result}: the cost per message of virtual against static dispatch, predictable and not, and the
four-stage pipeline's cost both ways.
\item Why is the \omterm{def:ll:cpu-microarchitecture:ubench}{microbenchmark} of a virtual call alone the wrong measure?
\item When is \omterm{def:ll:cpp-for-latency-ii-compile-time:dyn}{dynamic dispatch} the right choice in a trading system?
\item What does a concept add to a template interface?
\item What is \omterm{def:ll:cpp-for-latency-ii-compile-time:consteval}{compile-time evaluation} good for on the hot path?
\item Why keep error reporting out of line?
\item Why check the assembly rather than trust \texttt{always\_inline}?
\item In one sentence: what should the compiler know when it compiles the hot path?
\end{enumerate}
\end{problem}

\section{Interview questions}

\begin{interviewq}[$\star$ \iqroles{developer}]\label{iq:ll:cpp-for-latency-ii-compile-time:1}
How does a virtual function call work, and what does it cost?
\end{interviewq}

\begin{interviewq}[$\star\star$ \iqroles{developer}]\label{iq:ll:cpp-for-latency-ii-compile-time:2}
Explain the \omterm{def:ll:cpp-for-latency-ii-compile-time:static}{curiously recurring template pattern} and when you would use it.
\end{interviewq}

\begin{interviewq}[$\star\star$ \iqroles{developer}]\label{iq:ll:cpp-for-latency-ii-compile-time:3}
What is the difference between \texttt{constexpr} and \texttt{consteval}?
\end{interviewq}

\begin{interviewq}[$\star\star$ \iqroles{developer}]\label{iq:ll:cpp-for-latency-ii-compile-time:4}
Why can \omterm{def:ll:cpp-for-latency-ii-compile-time:inline}{inlining} make code faster by more than the cost of the call it removes? When can it make code slower?
\end{interviewq}

\begin{interviewq}[$\star\star$ \iqroles{developer}]\label{iq:ll:cpp-for-latency-ii-compile-time:5}
\texttt{std::variant} or a class hierarchy for message types: which would you choose on a hot path, and why?
\end{interviewq}

\begin{interviewq}[$\star\star\star$ \iqroles{developer}]\label{iq:ll:cpp-for-latency-ii-compile-time:6}
Design a trading engine whose strategy is chosen at start-up but whose hot path has no indirect calls.
\end{interviewq}
